\documentclass[11pt]{article}
\usepackage[margin=1in]{geometry}
\usepackage{amsmath,amssymb,amsthm}
\usepackage[T1]{fontenc}
\usepackage{lmodern}
\usepackage{microtype}
\usepackage[hidelinks]{hyperref}
\newtheorem{theorem}{Theorem}
\newtheorem{lemma}{Lemma}
\newtheorem{remark}{Remark}
\newcommand{\R}{\mathbb R}
\newcommand{\C}{\mathbb C}
\title{No open exact-root cover for Vassiliev's quadratic family}
\author{Alec Kriebel\\
\small Independent researcher\\
\small ORCID: \href{https://orcid.org/0009-0001-9320-500X}{0009-0001-9320-500X}}
\date{13 October 2026}
\begin{document}
\maketitle
\begin{abstract}
Problem~2A in Vassiliev's 2015 problem list asks for the minimum number of
open coefficient sets supporting continuous exact real-root selections for
a six-parameter quadratic system. For the full coefficient space printed
there, no such open cover exists, of any cardinality. Every neighborhood
of the zero coefficients contains the complex square-root obstruction.
We give an elementary proof and a translated family of obstructed
coefficients. This is an explicit diagnosis of the literal formulation
using classical topology and the source's own square-root slice.
Historical firstness is not established. No result is claimed for a
modified finite-genus question or the separate approximate-root problem.
\end{abstract}

\section{The printed question and its answer}
For
\(p=(a_1,a_2,a_0,b_1,b_2,b_0)\in\R^6\), put
\begin{align}
 F_p(x,y)&=x^2-y^2+a_1x+a_2y+a_0,\label{eq:f}\\
 G_p(x,y)&=xy+b_1x+b_2y+b_0.\label{eq:g}
\end{align}
An admissible set is a Euclidean-open \(U\subset\R^6\) with a continuous
map \(s:U\to\R^2\) satisfying \(F_p(s(p))=G_p(s(p))=0\) for every \(p\in U\).
Vassiliev's Problem~2A asks for a minimum number of such sets covering
all of \(\R^6\)~\cite[Section~2, p.~205]{V15}. Repeated-root coefficients
are included. A selection chooses an actual point, irrespective of
algebraic multiplicity.

\begin{theorem}\label{thm:main}
No open neighborhood of \(0\in\R^6\) is admissible.
Consequently, no admissible open cover of \(\R^6\) exists, even if
arbitrarily many members are allowed.
\end{theorem}
Thus the admissible-cover class is empty. If a finite-cover number is
extended by assigning \(+\infty\) when no finite admissible cover exists,
its value here is \(+\infty\); this notation does not assert the existence
of an infinite admissible cover.

\section{An elementary local proof}
\begin{lemma}\label{lem:square}
For every \(r>0\), there is no continuous \(S:D_r\to\C\), where
\(D_r=\{w\in\C:|w|<r\}\), satisfying \(S(w)^2=w\).
\end{lemma}
\begin{proof}
Consider the continuous path
\[
 \gamma(t)=
 \begin{cases}
 (1-t)+it,&0\leq t\leq1,\\
 (1-t)+i(2-t),&1\leq t\leq2.
 \end{cases}
\]
Its two formulas agree at \(t=1\), and
\[
 |\gamma(t)|^2=
 \begin{cases}
 2(t-\tfrac12)^2+\tfrac12,&0\leq t\leq1,\\
 2(t-\tfrac32)^2+\tfrac12,&1\leq t\leq2.
 \end{cases}
\]
Hence \(1/2\leq|\gamma(t)|^2\leq1\), with
\(\gamma(0)=1\) and \(\gamma(2)=-1\).
Choose \(c>0\) with \(c^2<r\). The path
\(w(t)=c^2\gamma(t)^2\) lies in \(D_r\setminus\{0\}\) and has
\(w(0)=w(2)=c^2\). If \(S\) existed, the continuous function
\[
 q(t)=\frac{S(w(t))}{c\gamma(t)}
\]
would satisfy \(q(t)^2=1\). Connectedness of \([0,2]\) makes \(q\)
constant. But its numerator agrees at the endpoints and its nonzero
denominator changes sign, so \(q(2)=-q(0)\), a contradiction.
\end{proof}

\begin{proof}[Proof of Theorem~\ref{thm:main}]
Restrict coefficients by the continuous embedding
\[
 \iota(u+iv)=(0,0,-u,0,0,-v/2).
\]
Equations~\eqref{eq:f}--\eqref{eq:g} then become
\(x^2-y^2=u\) and \(2xy=v\), so their full real solution set is
exactly the complex roots \(z=x+iy\) of \(z^2=u+iv\).
If \(U\) were an admissible neighborhood of zero, \(\iota^{-1}(U)\)
would be an open neighborhood of zero in \(\C\) and would contain
some \(D_r\). Composing its exact selection with \(\iota\), and
identifying \(\R^2\) with \(\C\), would contradict
Lemma~\ref{lem:square}. Every open cover has a member containing zero,
so its cardinality cannot evade this obstruction.
\end{proof}

\begin{remark}[Translated obstructions]\label{rem:translated}
For any \(h,k\in\R\), define
\[
 p_{h,k}=(-2h,\,2k,\,h^2-k^2,\,-k,\,-h,\,hk).
\]
No open neighborhood of \(p_{h,k}\) is admissible. Indeed, replacing
the constant coefficients by \(a_0=h^2-k^2-u\) and
\(b_0=hk-v/2\), while fixing the linear coefficients, gives
\[
 (x-h)^2-(y-k)^2=u,\qquad 2(x-h)(y-k)=v.
\]
The same argument applies to \(z=(x-h)+i(y-k)\).
In particular \(p_{1,0}=(-2,0,1,0,-1,0)\neq0\) is obstructed.
Removing just the zero coefficient vector therefore does not yield
an admissible open cover of the remaining coefficient space.
\end{remark}

\section{Provenance and the limits of the answer}
The final journal statement~\cite{V15} and its arXiv version
1504.01997v1 both ask for all of \(\R^6\), ordinary open domains, and
continuous exact roots. The source itself identifies the complex-square
slice. Its finite lower-bound hint does not remove zero.
For \(z^2=w\) over \(\C\setminus\{0\}\), two slit domains support branches
and one global branch is impossible. Including zero prevents even a
local exact branch. This is the domain distinction responsible for
Theorem~\ref{thm:main}; no intention to alter the source is attributed
to its author.

The strict-cover framework and the obstruction mechanism are classical.
Smale's definition uses open sets with strict right inverses, while his
root-selection Theorem~A removes the discriminant~\cite[pp.~83--85]{S87}.
Pave\v{s}i\'c explicitly distinguishes ordinary open strict sections
from filtered domains and homotopy sections, and warns that failure of
local sectionability defeats open section covers
\cite[Section~5, pp.~75--76]{P18}\cite[Section~2, pp.~109,111]{P19}.
We credit that machinery and Vassiliev's slice; no new topological
method is claimed.

A bounded primary-body priority audit did not locate an earlier explicit
diagnosis of this exact six-coefficient formulation. It cannot establish
novelty or historical firstness. Important coverage limits include
unread Schwarz originals/translation and Vassiliev's 1996 article;
the complete 2011 author manuscripts were read without comparison
to their final journal versions. Other version and selective-read
limits are recorded in the accompanying priority audit. This note
records the explicit literal answer, without a priority guarantee.

Deleting an appropriate coefficient locus, changing admissible domains,
or using homotopy sections would pose different questions.
Theorem~\ref{thm:main} does not solve any unspecified repaired finite-genus
problem. It also does not solve approximate Problem~2B in~\cite{V15}:
on the square slice, for fixed \(\varepsilon>0\), the value zero is
within \(\varepsilon\) of every root whenever \(|w|<\varepsilon^2\).

\section*{Verification and disclosure}
The supplementary package provides exact rational polynomial checks
of the slice identities, the full Jacobian, the path's universal
quadratic norm bounds, and adversarial input controls. The checks
remain active with Python optimization. They do not prove continuity
or the topological nonexistence statement; those are established by
the written argument above. Finite numerical sampling is not used
as a proof of the path bound.

AI tools were used extensively in developing, drafting, sourcing,
reproducing, and adversarially auditing this work. This preprint is
unrefereed and has not undergone conventional human peer review.
The author is responsible for the claims and their stated limits.

\begin{thebibliography}{9}
\bibitem{V15}
V.~A.~Vassiliev,
\emph{A Few Problems on Monodromy and Discriminants},
Arnold Math.\ J.\ \textbf{1} (2015), 201--209.
\href{https://doi.org/10.1007/s40598-015-0011-9}{doi:10.1007/s40598-015-0011-9}.
arXiv:1504.01997v1.

\bibitem{S87}
S.~Smale,
\emph{On the Topology of Algorithms, I},
J.\ Complexity \textbf{3} (1987), 81--89.
\href{https://doi.org/10.1016/0885-064X(87)90021-5}{doi:10.1016/0885-064X(87)90021-5}.

\bibitem{P18}
P.~Pave\v{s}i\'c,
\emph{A Topologist's View of Kinematic Maps and Manipulation Complexity},
Contemp.\ Math.\ \textbf{702} (2018), 61--83.
\href{https://doi.org/10.1090/conm/702/14107}{doi:10.1090/conm/702/14107}.

\bibitem{P19}
P.~Pave\v{s}i\'c,
\emph{Topological Complexity of a Map},
Homology Homotopy Appl.\ \textbf{21}(2) (2019), 107--130.
\href{https://doi.org/10.4310/HHA.2019.v21.n2.a7}{doi:10.4310/HHA.2019.v21.n2.a7}.
The relevant final-body passages were checked separately from the
later corrected arXiv:1809.09021v3.
\end{thebibliography}
\end{document}

